\chapter{Tools e Tecnologie}\label{chapter:Tools-Tecnologie}

Esploriamo ora i principali strumenti e le tecnologie di OSINT (Open Source Intelligence) che possono essere utilizzati per monitorare e analizzare le attività del dark web. Questi strumenti permettono alle aziende e alle organizzazioni di individuare le proprie vulnerabilità, raccogliendo informazioni sui dati esposti e sulle potenziali minacce. L'individuazione dei punti deboli e dei possibili punti deboli, i quali possono favorire un attacco, è fondamentale per rafforzare le difese di sicurezza informatica e prevenire attacchi mirati. Di seguito, vedremo i principali tools utilizzati per raccogliere, analizzare e proteggere le infrastrutture aziendali nel contesto dell'OSINT e del dark web.

Osintframework.com , un'ampia directory di strumenti e risorse OSINT gratuiti online ospitati sulla piattaforma per sviluppatori GitHub. Sia gli hacker che i professionisti della sicurezza informatica possono utilizzare questo elenco come punto di partenza per approfondire le funzionalità specifiche che desiderano in uno strumento OSINT.

Maltego , una soluzione di data mining in tempo reale per piattaforme Windows, Mac e Linux che fornisce rappresentazioni grafiche di modelli di dati e connessioni. Con la sua capacità di profilare e tracciare le attività online degli individui, questo strumento è utile sia ai professionisti della sicurezza informatica che agli attori delle minacce.

Spiderfoot  – Uno strumento di integrazione dell'origine dati per informazioni come indirizzi e-mail, numeri di telefono, indirizzi IP, sottodomini e altro ancora. Gli hacker etici possono utilizzare questa risorsa per indagare sulle informazioni disponibili al pubblico che potrebbero rappresentare una minaccia per un'organizzazione o un individuo.

Shodan , un motore di ricerca per dispositivi connessi a internet in grado anche di fornire informazioni sui metadati e sulle porte aperte. Poiché questo strumento è in grado di identificare le vulnerabilità di sicurezza per milioni di dispositivi, può essere utile sia per i professionisti della sicurezza informatica che per i criminali informatici.

Metasploit , uno strumento di test di penetrazione in grado di identificare le vulnerabilità della sicurezza nelle reti, nei sistemi e nelle applicazioni. Sia i professionisti della sicurezza informatica che gli hacker trovano valore in questo strumento perché può rivelare i punti deboli specifici che possono consentire un attacco informatico di successo.



1. Maltego
Descrizione: Maltego è uno strumento di visualizzazione e analisi dei dati che permette di identificare le relazioni tra informazioni come persone, aziende, indirizzi IP, domini, e altro. È particolarmente utile per tracciare le connessioni e mappare le reti di attori minacciosi.
Utilità: Nel contesto del dark web, Maltego può essere utilizzato per mappare le connessioni tra utenti sospetti e attività di cybercrime, facilitando l'identificazione di possibili minacce per l'organizzazione.
2. Shodan
Descrizione: Shodan è un motore di ricerca per dispositivi connessi a Internet e può identificare server, router, webcam, e altri dispositivi esposti pubblicamente. È utilizzato per identificare risorse non sicure o esposte.
Utilità: Permette di scoprire se i dispositivi dell'organizzazione sono visibili su Internet e se sono esposti a vulnerabilità note. Questo è cruciale per prevenire attacchi da parte di hacker che utilizzano il dark web per trovare bersagli facili.
3. SpiderFoot
Descrizione: SpiderFoot è uno strumento di OSINT automatizzato che raccoglie informazioni da una vasta gamma di fonti, inclusi il dark web, il surface web, e diverse API di terze parti. È ideale per la scoperta automatizzata di informazioni sensibili.
Utilità: Può essere utilizzato per scansionare il dark web alla ricerca di informazioni sensibili riguardanti l'organizzazione, come credenziali compromesse, menzioni in forum di hacker, o discussioni su potenziali vulnerabilità.
4. Censys
Descrizione: Censys è simile a Shodan ma fornisce analisi più dettagliate sugli asset connessi a Internet. Offre anche la possibilità di monitorare e ottenere avvisi sulle esposizioni dei dati.
Utilità: Per un'azienda, monitorare costantemente i propri asset connessi a Internet per esporre potenziali punti deboli è fondamentale per prevenire accessi non autorizzati che potrebbero essere sfruttati dai criminali del dark web.
5. Dark Web Monitoring Tools (ad esempio, DarkOwl, Recorded Future)
Descrizione: Questi strumenti specializzati consentono alle organizzazioni di monitorare il dark web alla ricerca di dati aziendali compromessi, come credenziali, documenti sensibili, o informazioni su possibili attacchi imminenti.
Utilità: Forniscono un monitoraggio continuo delle menzioni di un'organizzazione nel dark web, aiutando a rilevare tempestivamente eventuali fughe di dati o attività malevole pianificate contro l'organizzazione.
6. Have I Been Pwned
Descrizione: È un servizio che permette di verificare se le credenziali aziendali o personali sono state compromesse in una violazione di dati pubblica.
Utilità: Essenziale per le organizzazioni per monitorare regolarmente se i loro dipendenti sono coinvolti in breach di dati noti e per rispondere rapidamente modificando le password e adottando misure di sicurezza aggiuntive.
7. Threat Intelligence Platforms (TIPs)
Descrizione: Strumenti come MISP (Malware Information Sharing Platform) consentono di raccogliere, organizzare e condividere informazioni su minacce note e emergenti.
Utilità: Questi strumenti aiutano a collegare i dati OSINT raccolti con minacce conosciute e a stabilire contromisure appropriate. Essere proattivi nel monitoraggio delle minacce consente di migliorare la sicurezza complessiva dell'organizzazione.
8. Social Media Intelligence Tools (ad esempio, Social Searcher, Mentionlytics)
Descrizione: Questi strumenti permettono di monitorare le menzioni dell'azienda e i suoi dipendenti sui social media, che spesso sono usati da hacker per raccogliere informazioni.
Utilità: È possibile identificare eventuali campagne di phishing, false informazioni, o attacchi di social engineering pianificati.
9. VPN e Sicurezza di Navigazione Tor
Descrizione: Utilizzare VPN e il browser Tor è essenziale per mantenere la sicurezza e l'anonimato durante il monitoraggio del dark web.
Utilità: Minimizzare il rischio di rilevamento da parte degli attori malevoli durante le operazioni di intelligence sul dark web è fondamentale per la sicurezza dell'investigatore.
10. Elastic Stack (ELK)
Descrizione: Elastic Stack, composto da Elasticsearch, Logstash, e Kibana, è una suite di strumenti per il monitoraggio e l'analisi dei dati.
Utilità: Può essere utilizzato per raccogliere e analizzare i dati di log e monitorare i segnali di compromissione (Indicators of Compromise, IoC) correlati alle minacce del dark web.


Testo di presentazione per i tools e le tecnologie

\section{OSINT Framework}\label{sec:cap_sec_tool1}

Prima Sezione

\subsection{Caratteristiche e Prova di Utilizzo}\label{subsec:es_subsec_tool1}

Prima Sottosezione

\section{Ahmia Search Engine}\label{sec:cap_sec_tool2}

Seconda Sezione

\subsection{Caratteristiche e Prova di Utilizzo}\label{subsec:es_subsec_tool2}

Seconda Sottosezione

\section{Shodan}\label{sec:cap_sec_tool3}

Terza Sezione

\subsection{Caratteristiche e Prova di Utilizzo}\label{subsec:es_subsec_tool3}

Terza Sottosezione